% الجزء: sets-functions-relations
% الفصل: relations
% القسم: reflections
%
\documentclass[../../../include/open-logic-section]{subfiles}

\begin{document}

\olfileid[ar]{sfr}{rel}{ref}
\olsection{تأملات فلسفية}

عرّفنا العلاقات في \olref[set]{sec} بمجموعات معينة. وهنا مسألة فلسفية
وجيزة: ما \emph{مؤدى} هذا التعريف؟ لا يكاد يصح أن نزعم أننا
\emph{اكتشفنا} به حقائق في الهوية الميتافيزيقية؛ كأن علاقة الترتيب
على $\Nat$ \emph{انكشف أنها} المجموعة
$R=  \Setabs{\tuple{n,m}}{n, m \in \Nat\text{ و }
n < m}$ المعرّفة في \olref[set]{sec}. وبيان ذلك من ثلاثة وجوه.

أما الأول، ففي \olref[set][pai]{wienerkuratowski} جعلنا $\tuple{a, b} =
\{\{a\}, \{a, b\}\}$. وكان لنا أن نجعل بدلًا منه $\lVert a,
b\rVert = \{\{b\}, \{a, b\}\} = \tuple{b,a}$. فإذا كان $a \neq b$
كان $\tuple{a, b} \neq \lVert a,b\rVert$. ومع ذلك يصلح
$\lVert a,b\rVert$ لتعريف الزوج المرتب كما يصلح $\tuple{a,b}$،
ويؤدي كل منهما الغرض نفسه. فصار لدينا شيئان متكافئان في الصلاحية
لأن يعد أحدهما «نفس» علاقة الترتيب على الطبيعيّات:
\begin{align*}
		R &= \Setabs{\tuple{n,m}}{n, m \in \Nat \text{ و }n < m}\\
		S &= \Setabs{\lVert n,m\rVert}{n, m \in \Nat \text{ و }n < m}.
\end{align*}
والامتدادية تقضي بأن $R \neq S$؛ فلا يمكن أن يكونا \emph{معًا}
عين علاقة الترتيب على~$\Nat$. أما تخصيص $R$ دون $S$، أو بالعكس،
بأنها \emph{هي} العلاقة في نفس الأمر، فتحكم محض يوقع في الحرج.
وهذا مثال بالغ البساطة لحجة على الاختزالية بالمجموعات اشتهرت عن
بيناسيراف في \citeyear{Benacerraf1965}، وسنعاودها مرارًا.

وأما الثاني، فلو جعلنا \emph{كل} علاقة عين مجموعة، لكانت علاقة
الانتماء نفسها $ \in$ مجموعة معينة، ولا بد أن تكون
$\Setabs{\tuple{x,y}}{x \in y}$. أفتوجد هذه المجموعة؟ ليست دعوى
وجودها هينة بعد مفارقة راسل. بل
\oliflabeldef{cumul:::part}{إن نظرية المجموعات المبنية في
\olref[cumul][][]{part} \emph{تنفي} وجودها.\footnote{ولنستبق البيان:
افترض بالخلف وجود $I = \Setabs{\tuple{x,y}}{x \in y}$. فيكون
$\bigcup \bigcup I$ المجموعة الشاملة، خلافًا لـ
\olref[sfr][z][sep]{thm:NoUniversalSet}.}}{يمكن إقامة نظرية مجموعات
بديهية محكمة تنفي وجود هذه المجموعة.}
فمهما جاز أخذ بعض العلاقات مجموعات، فلا بد للانتماء من حكم خاص.

وأما الثالث، فحين جعلنا العلاقات مجموعات أجزنا كتابة $Rxy$ مكان
$\tuple{x,y} \in R$. وهذا سائغ ما دامت علاقة الانتماء «$\in$»
مأخوذة \emph{محمولًا}. فإن جعلنا «$\in$» اسمًا لمجموعة مخصوصة، لم
يبق في «$\tuple{x,y} \in R$» إلا ثلاثة حدود مفردة دالة على مجموعات:
«$\tuple{x,y}$» و«$\in$» و«$R$». ومجرد سرد الأسماء لا يعبر عن قضية،
كما لا تعبر عنها العبارة المختلة «الكوب حامل الأقلام الطاولة».
فلا بد إذن من استثناء علاقة الانتماء، وإن جاز عد غيرها مجموعات.
ويجتمع في هذا الوجه تبسيط لمفارقة مفهوم \emph{الحصان} عند فريغه
واعتراض مشهور وجّهه فيتغنشتاين إلى راسل.

فالحاصل أنه لا \emph{خطأ} في قولنا إن العلاقات على الأعداد مجموعات،
وإنما يجب حمل القول على مراده. لسنا نثبت به حقيقة ميتافيزيقية في
الهوية؛ وإنما نقرر أن بعض السياقات يجيز \emph{معاملة} علاقات معينة
معاملة مجموعات معينة، وسنستعمل هذه الإجازة فيما يأتي.

\end{document}
